\documentclass[../NDCNNAI_main.tex]{subfiles}
\graphicspath{{\subfix{../fig/}}}

\begin{document}
\subsection{Три взгляда на универсальную аппроксимацию глубокими свёрточными сетями}
\subsubsection{Модель Яроцкого: приближение инвариантных и эквивариантых отображений}
Свёрточные сети устроены так, что их линейная часть (свёртка с \emph{разделенным весом})
естественно согласована с действием группы сдвигов.
Поэтому вопрос об \emph{универсальной аппроксимации} для CNN имеет смысл только вместе с уточнением:
\emph{какой класс отображений} мы хотим аппроксимировать --- эквивариантные операторы, инвариантные функционалы,
или произвольные непрерывные отображения. Яроцкий предлагает строгую <<континуальную>> постановку, где CNN рассматриваются как аппроксиматоры отображений
между бесконечномерными пространствами сигналов, а конечные архитектуры возникают через дискретизацию и усечение
(окно) \cite{Yarotsky2018}.

Пусть $S_t$ обозначает оператор сдвига входного сигнала/изображения на $t$ (в непрерывной постановке ниже это будет $R_\gamma$).
\begin{definition}
Отображение $\mathcal{F}$ называется
\begin{itemize}
  \item \emph{сдвиг-эквивариантным}, если $\mathcal{F}(S_t x)=S_t \mathcal{F}(x)$ для всех $t$;
  \item \emph{сдвиг-инвариантным}, если $\mathcal{F}(S_t x)=\mathcal{F}(x)$ для всех $t$.
\end{itemize}
\end{definition}

\textit{Интуитивно}: эквивариантность означает \emph{коммутирование} с действием группы сдвигов; инвариантность означает, что выход вообще не меняется при сдвиге.
Свёрточные слои с \emph{разделением весов} (weight sharing) естественно навязывают эквивариантность, а \emph{пулинг} и/или \emph{глобальная агрегация} обычно используются как механизм получения инвариантности.

Для дискретного отображения признаков $z\colon\mathbb{Z}^\nu\to\mathbb{R}^m$ и шага (stride) $s\in\mathbb{N}$ \emph{локальный пулинг} (downsampling) строит более грубую отображение на $s\mathbb{Z}^\nu$:
\[
(\mathrm{Pool}_s z)(\gamma)=\mathrm{pool}\big(\{z(s\gamma+\theta):\theta\in\mathcal{B}\}\big),\qquad \gamma\in\mathbb{Z}^\nu,
\]
где $\mathcal{B}$ --- локальное окно (например, $\mathcal{B}=Z_{L_{\mathrm{rf}}}$), а $\mathrm{pool}$ --- max/avg и т.п.
В модели Яроцкого используется особенно <<чистый>> вариант пулинга: \emph{децимация} (subsampling)
\[
(\mathrm{Dec}_s z)(\gamma)=z(s\gamma),\qquad \gamma\in\mathbb{Z}^\nu,
\]
то есть оставляем каждую $s$-ю пространственную позицию (stride-$s$ downsampling) \cite{Yarotsky2018}.
Наконец, \emph{глобальная агрегация} (global aggregation/pooling) <<схлопывает>> пространство, например
\[
A_{\mathrm{sum}}(z)=\sum_{\gamma\in\mathbb{Z}^\nu} z(\gamma),\qquad
A_{\mathrm{avg}}(z)=\frac{1}{|\Omega|}\sum_{\gamma\in\Omega} z(\gamma)
\quad (\Omega \text{ --- конечное окно}).
\]
Если $z$ сдвиг-эквивариантна, то $A(z)$ становится сдвиг-инвариантной.

Обсудим постановку Яроцкого для <<сигналов бесконечного протяжения>> \cite{Yarotsky2018}.
Фиксируем размерность пространства $\nu\in\mathbb{N}$ и числа каналов $d_V,d_U\in\mathbb{N}$.
Рассмотрим гильбертовы пространства
\[
V\coloneq L^2(\mathbb{R}^\nu,\mathbb{R}^{d_V}),\qquad
U\coloneq L^2(\mathbb{R}^\nu,\mathbb{R}^{d_U}),
\]
со стандартной нормой $\|\cdot\|_V,\|\cdot\|_U$.
Действие группы сдвигов $\mathbb{R}^\nu$ на $V$ (и аналогично на $U$) задаётся операторами
\[
(S_\gamma \Phi)(\theta)=\Phi(\theta-\gamma),\qquad \gamma,\theta\in\mathbb{R}^\nu,\ \Phi\in V.
\]

Отображение $f\colon V\to U$ называется \emph{$\mathbb{R}^\nu$-эквивариантным}, если
\[
f(S_\gamma\Phi)=S_\gamma f(\Phi)\qquad \forall \gamma\in\mathbb{R}^\nu,\ \forall \Phi\in V.
\]  
Отображение $f\colon V\to U$ \emph{непрерывно по норме}, если из $\|\Phi_n-\Phi\|_V\to 0$
следует $\|f(\Phi_n)-f(\Phi)\|_U\to 0$.

Ключевая идея подхода: аппроксимировать непрерывное $f\colon V\to U$ конечновычислимыми CNN, которые работают
на \emph{дискретизациях} сигнала на всё более тонких сетках и в всё более широких конечных окнах.

Для $\lambda>0$ пусть $V_\lambda\subset V$ --- подпространство сигналов, постоянных на кубах
\[
Q^{(\lambda)}_{\mathbf{k}}=\prod_{s=1}^\nu\Big[\big(k_s-\tfrac12\big)\lambda,\ \big(k_s+\tfrac12\big)\lambda\Big],
\qquad \mathbf{k}=(k_1,\dots,k_\nu)\in\mathbb{Z}^\nu.
\]
Обозначим через $P_\lambda:V\to V_\lambda$ ортогональный проектор:
\[
(P_\lambda\Phi)(\gamma)=\frac{1}{\lambda^\nu}\int_{Q^{(\lambda)}_{\mathbf{k}}}\Phi(\theta)\,d\theta,
\qquad \text{где } Q^{(\lambda)}_{\mathbf{k}}\ni \gamma.
\]
Тогда $V_\lambda\simeq L^2\big((\lambda\mathbb{Z})^\nu,\mathbb{R}^{d_V}\big)$ (с естественным скалярным произведением).

Для целого $L\ge 0$ положим
\[
Z_L=\{\mathbf{k}\in\mathbb{Z}^\nu:\|\mathbf{k}\|_\infty\le L\}.
\]
Для $\Lambda\ge 0$ (пространственный радиус / cut-off) определим подпространство дискретных сигналов,
обнуляющихся вне окна:
\[
V_{\lambda,\Lambda}
=\Big\{\Phi:(\lambda\mathbb{Z})^\nu\to\mathbb{R}^{d_V}:\ \Phi(\lambda\mathbf{k})=0\ \text{если }\mathbf{k}\notin Z_{\lfloor \Lambda/\lambda\rfloor}\Big\}
\simeq L^2\big(\lambda Z_{\lfloor \Lambda/\lambda\rfloor},\mathbb{R}^{d_V}\big).
\]
Аналогично определяются $U_{\lambda,\Lambda}\subset U$ и ортогональные проекторы $P_{\lambda,\Lambda}\colon V\to V_{\lambda,\Lambda}$,
$P_{\lambda,\Lambda}\colon U\to U_{\lambda,\Lambda}$.

Эвристика (важна для <<предела CNN>>): $\lambda\downarrow 0$ фиксирует всё более тонкие детали сигнала,
а $\Lambda\uparrow\infty$ делает ошибку усечения хвостов пренебрежимо малой на фиксированных компактах. Поэтому <<предел CNN>> формализуется через последовательность моделей на всё более тонких решётках и всё больших окнах.

Нам понадобится стандартная <<компактная>> версия сильной сходимости проекторов.

\begin{lemma}[Равномерная аппроксимация проекторами на компакте]
\label{lem:proj_approx}
Пусть $K\subset V$ --- компакт, и $\varepsilon>0$.
Тогда существуют $\lambda_0>0$ и $\Lambda_0>0$ такие, что при $\lambda\le\lambda_0$ и $\Lambda\ge\Lambda_0$ выполнено
\[
\sup_{\Phi\in K}\|P_{\lambda,\Lambda}\Phi-\Phi\|_V<\varepsilon.
\]
Аналогичное утверждение верно в $U$.  
\end{lemma}

\begin{proof}
В \cite{Yarotsky2018} показано, что $P_{\lambda,\Lambda}\Phi\to \Phi$ в $V$ при $\lambda\to 0$, $\Lambda\to\infty$
для каждого фиксированного $\Phi$ (сильная сходимость проекторов).
Переходим от поточечной сходимости к равномерной на компакте с помощью $\varepsilon$–сетей.

Так как $K$ компактен, существует конечное $\varepsilon/3$–покрытие: $K\subset\bigcup_{j=1}^N B(\Phi_j,\varepsilon/3)$.
Для каждой $\Phi_j$ найдутся $\lambda_j,\Lambda_j$ такие, что при $\lambda\le\lambda_j$ и $\Lambda\ge\Lambda_j$ выполнено
$\|P_{\lambda,\Lambda}\Phi_j-\Phi_j\|_V<\varepsilon/3$.
Положим $\lambda_0\coloneq \min_j \lambda_j$ и $\Lambda_0\coloneq \max_j\Lambda_j$.

Возьмём произвольное $\Phi\in K$ и $j$ такое, что $\|\Phi-\Phi_j\|_V<\varepsilon/3$.
Так как $P_{\lambda,\Lambda}$ --- ортогональный проектор, $\|P_{\lambda,\Lambda}\|\le 1$. Тогда
\[
\|P_{\lambda,\Lambda}\Phi-\Phi\|_V
\le \|P_{\lambda,\Lambda}(\Phi-\Phi_j)\|_V + \|P_{\lambda,\Lambda}\Phi_j-\Phi_j\|_V + \|\Phi_j-\Phi\|_V
< \varepsilon.
\]
\end{proof}

Фиксируем \emph{область поля восприятия} (receptive field) как кубическое множество
\[
Z_{L_{\mathrm{rf}}}\subset\mathbb{Z}^\nu,\qquad L_{\mathrm{rf}}\in\mathbb{N}, \qquad Z_{L_\mathrm{rf}} = \{\mathbf{k}\in\mathbb{Z}^\nu:\|\mathbf{k}\|_\infty\le L_\mathrm{rf}\}
\]
и считаем $L_{\mathrm{rf}}$ фиксированным, тогда как $\lambda$ и $\Lambda$ могут меняться (детализация и окно).

\begin{definition}
  \emph{Базовая свёрточная сеть} (basic convnet) --- это отображение $\hat f\colon V\to U$, определённое композицией
\begin{equation}
  \hat f\colon \ V \xrightarrow{\,P_{\lambda,\Lambda+(T-1)\lambda L_{\mathrm{rf}}}\,}
V_{\lambda,\Lambda+(T-1)\lambda L_{\mathrm{rf}}}(=W_1)
\xrightarrow{\ \hat f_1\ } W_2 \xrightarrow{\ \hat f_2\ }\cdots
\xrightarrow{\ \hat f_T\ } W_{T+1}(=U_{\lambda,\Lambda}),
\label{eq:basic_chain}  
\end{equation}

где промежуточные пространства имеют вид
\[
W_t=\Big\{\Phi:\lambda Z_{\lfloor \Lambda/\lambda\rfloor+(T-t)L_{\mathrm{rf}}}\to \mathbb{R}^{d_t}\Big\},
\quad t\le T,
\qquad
W_{T+1}=\Big\{\Phi:\lambda Z_{\lfloor \Lambda/\lambda\rfloor}\to \mathbb{R}^{d_U}\Big\},
\]
а слои задаются как <<линейная свёртка + нелинейность>>:
\[
\big(\hat f_t(\Phi)\big)_{\gamma,n}
=
\sigma\!\left(\sum_{\theta\in Z_{L_{\mathrm{rf}}}}\sum_{k=1}^{d_t}
w^{(t)}_{n\theta k}\,\Phi_{\gamma+\theta,k}+h^{(t)}_n\right),
\qquad
\gamma\in Z_{\lfloor \Lambda/\lambda\rfloor+(T-t-1)L_{\mathrm{rf}}},
\]
и на последнем слое $t=T$ нелинейность опускается:
\[
\big(\hat f_T(\Phi)\big)_{\gamma,n}
=
\sum_{k=1}^{d_T} w^{(T)}_{nk}\,\Phi_{\gamma,k}+h^{(T)}_n,
\qquad
\gamma\in Z_{\lfloor \Lambda/\lambda\rfloor},\ n=1,\dots,d_U.
\]
Здесь $\sigma$ --- фиксированная поэлементная активация (например, ReLU), а коэффициенты $w^{(t)}_{n\theta k}$ и $h^{(t)}_n$
являются параметрами сети.
\end{definition}

Из-за разделения весов и того, что свёртка локальна и одинакова во всех узлах решётки, модель \eqref{eq:basic_chain} коммутирует со сдвигами на решётке $\lambda\mathbb{Z}^\nu$
(с поправкой на граничные эффекты усечения), что будет ключевым в доказательстве необходимости теоремы ниже \cite{Yarotsky2018}.

Из-за бесконечномерности $V$ топология равномерной сходимости на компактах не удовлетворяет первой аксиоме счётности,
поэтому Яроцкий аккуратно фиксирует определение <<предельной точки>> через аппроксимацию \emph{на каждом компакте} \cite{Yarotsky2018}.

\begin{definition}
Пусть $V=L^2(\mathbb{R}^\nu,\mathbb{R}^{d_V})$ и $U=L^2(\mathbb{R}^\nu,\mathbb{R}^{d_U})$.
Говорят, что отображение $f:V\to U$ является \emph{предельной точкой семейства базовых свёрточных сетей},
если для любого фиксированного $L_{\mathrm{rf}}$, любого компакта $K\subset V$, любых $\varepsilon>0$, $\lambda_0>0$, $\Lambda_0>0$
существует базовая свёрточная сеть $\hat f$ с тем же $L_{\mathrm{rf}}$, шагом $\lambda\le \lambda_0$ и радиусом $\Lambda\ge \Lambda_0$ такая, что
\[
\sup_{\Phi\in K}\ \|\hat f(\Phi)-f(\Phi)\|_U<\varepsilon.
\]  
\end{definition}

\begin{theorem}[Характеризация предельных точек, теорема 3.1 \cite{Yarotsky2018}]
Отображение $f\colon V\to U$ является предельной точкой семейства базовых CNN
тогда и только тогда, когда $f$ (i) непрерывно по норме и (ii) сдвиг-эквивариантно.  
\end{theorem}

\begin{proof}
\textbf{(Необходимость).}
Пусть $f$ --- предельная точка базовых CNN.

\underline{Шаг 1: непрерывность $f$.}
Пусть $\Phi_n\to\Phi$ в $V$. Множество
\(
K\coloneq \{\Phi\}\cup\{\Phi_n:\ n\in\mathbb{N}\}
\)
компактно (в метрическом пространстве равносильны компактность и секвенциальная компактность).
Возьмём $\varepsilon>0$, существует базовая CNN $\widehat f$ такая, что
\begin{equation}
\sup_{\Psi\in K}\|f(\Psi)-\widehat f(\Psi)\|_U<\varepsilon/3.
\label{eq:cont1}  
\end{equation}
Каждая конкретная $\widehat f$ непрерывна как композиция непрерывных отображений
(проектор + конечное число конечномерных операций + вложение обратно в $U$),
поэтому существует $N$ такое, что при $n\ge N$ выполнено
\begin{equation}
\|\widehat f(\Phi_n)-\widehat f(\Phi)\|_U<\varepsilon/3.
\label{eq:cont2}
\end{equation}
Тогда при $n\ge N$, по неравенству треугольника и \eqref{eq:cont1}--\eqref{eq:cont2},
\[
\|f(\Phi_n)-f(\Phi)\|_U
\le \|f(\Phi_n)-\widehat f(\Phi_n)\|_U
+\|\widehat f(\Phi_n)-\widehat f(\Phi)\|_U
+\|\widehat f(\Phi)-f(\Phi)\|_U
<\varepsilon.
\]
Значит, $f$ непрерывно.

\underline{Шаг 2: эквивариантность $f$.}
Нужно показать $f(R_\gamma\Phi)=R_\gamma f(\Phi)$ для всех $\gamma$ и $\Phi$.
Фиксируем $\Phi\in V$ и $\gamma\in\mathbb{R}^\nu$.

\emph{Локализация в большом окне.}
Для $M>0$ положим $D_M=[-M,M]^\nu$ и обозначим $P_{D_M}$ оператор умножения на индикатор $1_{D_M}$ в $U$
(ортогональный проектор на подпространство функций, с носителем в $D_M$).
Так как $P_{D_M}u\to u$ в $U$ при $M\to\infty$ для любого $u\in U$, достаточно доказать тождество
\begin{equation}
P_{D_M}f(R_\gamma\Phi)=R_\gamma P_{D_M}f(\Phi)\qquad \forall M>0.
\label{eq:eq_goal_M}  
\end{equation}

\emph{Аппроксимация непрерывного сдвига дискретным.}
Для данного $\lambda>0$ выберем $\gamma_\lambda\in (\lambda\mathbb{Z})^\nu$ так, что
$\|\gamma-\gamma_\lambda\|_\infty\le \lambda$.
Используем факты из оригинальной статьи:
\begin{itemize}
\item $R_{\gamma_\lambda}\to R_\gamma$ сильно в $U$ при $\lambda\to 0$;
\item из непрерывности $f$ следует: $f(R_{\gamma_\lambda}\Phi)\to f(R_\gamma\Phi)$ в $U$ при $\lambda\to 0$;
\item $P_{\lambda,M}\to P_{D_M}$ сильно в $U$ при $\lambda\to 0$ (здесь $P_{\lambda,M}$ --- дискретизация+усечение на окне $D_M$).
\end{itemize}
Следовательно, для любого $\varepsilon>0$ можно выбрать $\lambda>0$ настолько малым, что одновременно:
\begin{gather}
\|P_{D_M}f(R_\gamma\Phi)-P_{\lambda,M}f(R_\gamma\Phi)\|_U <\varepsilon,\label{eq:a1}\\
\|P_{\lambda,M}f(R_\gamma\Phi)-P_{\lambda,M}f(R_{\gamma_\lambda}\Phi)\|_U <\varepsilon,\label{eq:a2}\\
\|R_{\gamma_\lambda}P_{\lambda,M}f(\Phi)-R_\gamma P_{D_M}f(\Phi)\|_U <2\varepsilon.\label{eq:a3}
\end{gather}

\emph{Использование эквивариантности конкретной CNN на решётке.}
Рассмотрим компактное множество $K\coloneq \{\Phi,R_{\gamma_\lambda}\Phi\}$.
По определению предельной точки найдётся базовая CNN $\widehat f$ с данным $\lambda$ и достаточно большим $\Lambda$
такая, что
\begin{equation}
\sup_{\Psi\in K}\|f(\Psi)-\widehat f(\Psi)\|_U<\varepsilon.
\label{eq:a4}  
\end{equation}
Выбираем $\Lambda$ так, чтобы на области $D_M$ не было краевых эффектов при сдвиге,
и тогда для дискретного сдвига $\gamma_\lambda\in(\lambda\mathbb{Z})^\nu$ выполняется тождество
\begin{equation}
P_{\lambda,M}\widehat f(R_{\gamma_\lambda}\Phi)=R_{\gamma_\lambda}P_{\lambda,M}\widehat f(\Phi).
\label{eq:disc_equiv}
\end{equation}

\emph{Сборка оценки.}
Оценим норму разности в \eqref{eq:eq_goal_M}:
\begin{multline*}
\|P_{D_M}f(R_\gamma\Phi)-R_\gamma P_{D_M}f(\Phi)\|_U \le\\
\le
\underbrace{\|P_{D_M}f(R_\gamma\Phi)-P_{\lambda,M}f(R_\gamma\Phi)\|_U}_{<\varepsilon\ \text{по }\eqref{eq:a1}}
+
\underbrace{\|P_{\lambda,M}f(R_\gamma\Phi)-P_{\lambda,M}f(R_{\gamma_\lambda}\Phi)\|_U}_{<\varepsilon\ \text{по }\eqref{eq:a2}} + \\
+
\|P_{\lambda,M}f(R_{\gamma_\lambda}\Phi)-P_{\lambda,M}\widehat f(R_{\gamma_\lambda}\Phi)\|_U
+
\underbrace{\|P_{\lambda,M}\widehat f(R_{\gamma_\lambda}\Phi)-R_{\gamma_\lambda}P_{\lambda,M}\widehat f(\Phi)\|_U}_{=0\ \text{по }\eqref{eq:disc_equiv}} + \\
+
\|R_{\gamma_\lambda}P_{\lambda,M}\widehat f(\Phi)-R_{\gamma_\lambda}P_{\lambda,M}f(\Phi)\|_U
+
\underbrace{\|R_{\gamma_\lambda}P_{\lambda,M}f(\Phi)-R_\gamma P_{D_M}f(\Phi)\|_U}_{<2\varepsilon\ \text{по }\eqref{eq:a3}}.
\end{multline*}
Так как $\|P_{\lambda,M}\|\le 1$ и $\|R_{\gamma_\lambda}\|=1$, два <<средних>> члена оцениваются через \eqref{eq:a4} сверху на $\varepsilon$ каждый.
Итого получаем
\[
\|P_{D_M}f(R_\gamma\Phi)-R_\gamma P_{D_M}f(\Phi)\|_U < 6\varepsilon.
\]
Так как $\varepsilon$ произвольно, доказано \eqref{eq:eq_goal_M}. Переходом $M\to\infty$ получаем эквивариантность $f$.

\medskip
\textbf{(Достаточность).}
Пусть теперь $f\colon V\to U$ непрерывно по норме и эквивариантно.
Нужно: для любого компакта $K\subset V$ и $\varepsilon>0$ построить базовую CNN $\widehat f$, приближающую $f$ на $K$.

\underline{Шаг 1: замена $f$ на дискретизованно-усечённый суррогат.}
По лемме \ref{lem:proj_approx} (применённой к $K$ и $f(K)\subset U$, что компактно по непрерывности $f$)
можно выбрать $\lambda$ малым и $\Lambda$ большим так, что одновременно:
\[
\sup_{\Phi\in K}\|P_{\lambda,\Lambda}\Phi-\Phi\|_V<\delta,
\qquad
\sup_{\Psi\in f(K)}\|P_{\lambda,\Lambda}\Psi-\Psi\|_U<\varepsilon/3,
\]
и по равномерной непрерывности $f$ на компакте $K$ (теорема Кантора) при достаточно малом $\delta$ выполнено
\[
\sup_{\Phi\in K}\|f(P_{\lambda,\Lambda}\Phi)-f(\Phi)\|_U<\varepsilon/3.
\]
Тогда для приближения
\[
f_{\lambda,\Lambda}(\Phi)\coloneq P_{\lambda,\Lambda}\,f(P_{\lambda,\Lambda}\Phi)
\]
получаем
\begin{equation}
\sup_{\Phi\in K}\|f_{\lambda,\Lambda}(\Phi)-f(\Phi)\|_U<2\varepsilon/3.
\label{eq:surrogate}  
\end{equation}

\underline{Шаг 2: сведение к конечномерному <<ядру>> при $\Lambda=0$.}
Зафиксируем глубину $T$.
В \cite{Yarotsky2018} вводится функция $f_{\lambda,0,T}$, которая отвечает выходу на \emph{одной} пространственной ячейке
(вырожденное окно $\Lambda=0$) и зависит только от конечного входного окна радиуса $(T-1)\lambda L_{\mathrm{rf}}$.
Формально задаётся
\begin{equation}
f_{\lambda,0,T}(\Phi)
\coloneq 
P_{\lambda,0}\,f\!\big(P_{\lambda,(T-1)\lambda L_{\mathrm{rf}}}\Phi\big),
\qquad
f_{\lambda,0,T}:V_{\lambda,(T-1)\lambda L_{\mathrm{rf}}}\to U_{\lambda,0}.
\label{eq:coremap}  
\end{equation}
Это отображение непрерывно и действует между конечномерными пространствами.

\underline{Шаг 3: аппроксимация ядра базовой CNN (Лемма \ref{lem:approx_degen_window}).}
Ниже мы докажем Лемму \ref{lem:approx_degen_window}, которая говорит, что
любой непрерывный $f_{\lambda,0,T}$ можно аппроксимировать базовой CNN с теми же $\lambda,T$ и $\Lambda=0$.

\underline{Шаг 4: расширение ядра до произвольного окна (<<tile by translations>>).}
Эквивариантность $f$ позволяет восстановить полный выход на окне из локального ядра с помощью <<shift-and-sum>>
(смотри \cite[(3.28)]{Yarotsky2018}): для дискретных $\gamma\in Z_{\lfloor \Lambda/\lambda\rfloor}$
\begin{equation}
f_{\lambda,\Lambda,T}(\Phi)
\coloneq 
\sum_{\gamma\in Z_{\lfloor \Lambda/\lambda\rfloor}}
R_{\lambda\gamma}\, f_{\lambda,0,T}(R_{-\lambda\gamma}\Phi)
\quad \in U_{\lambda,\Lambda}.
\label{eq:shift_sum_target}  
\end{equation}
Так как сумма конечна, а сдвиги изометричны, то если $\widehat f_0$ аппроксимирует $f_{\lambda,0,T}$
на множествах $R_{-\lambda\gamma}K$, тогда сеть
\begin{equation}
\widehat f(\Phi)\coloneq \sum_{\gamma\in Z_{\lfloor \Lambda/\lambda\rfloor}}
R_{\lambda\gamma}\,\widehat f_0(R_{-\lambda\gamma}\Phi)
\label{eq:shift_sum_net}  
\end{equation}
аппроксимирует $f_{\lambda,\Lambda,T}$ на $K$, причём ошибка суммируется линейно и контролируется выбором точности ядра.
Реализуемость \eqref{eq:shift_sum_net} в классе базовых CNN обсуждается и формально доказывается в статье:
локальность (фикс.\ $L_{\mathrm{rf}}$) обеспечивает вычислимость каждого элемента расширения одним и тем же ядром,
а эквивариантность слоёв обеспечивает корректное <<копирование>> по всем позициям.

\underline{Шаг 5: итоговая оценка.}
Можно показать, что при достаточно большом $T$ функция $f_{\lambda,\Lambda,T}$ приближает $f_{\lambda,\Lambda}$
на $K$ (это контроль <<длины поля восприятия>> в континуальном масштабе).
Совмещая: \eqref{eq:surrogate} + аппроксимацию $f_{\lambda,\Lambda}$ через $f_{\lambda,\Lambda,T}$ +
аппроксимацию $f_{\lambda,\Lambda,T}$ сетью \eqref{eq:shift_sum_net},
получаем требуемое
$\sup_{\Phi\in K}\|\widehat f(\Phi)-f(\Phi)\|_U<\varepsilon$.  
\end{proof}

\begin{lemma}[Аппроксимация при вырожденном окне, лемма 3.1 \cite{Yarotsky2018}]
\label{lem:approx_degen_window}
Пусть фиксированы $\lambda>0$ и $T\in\mathbb{N}$, а выходное окно вырождено: $\Lambda=0$.
Тогда \emph{любое} непрерывное отображение
\[
g\colon V_{\lambda,(T-1)\lambda L_{\mathrm{rf}}}\to U_{\lambda,0}
\]
можно аппроксимировать (равномерно на компактах) базовыми CNN с теми же $\lambda$, $T$ и $\Lambda=0$
(разрешая увеличивать числа каналов $d_t$).
\end{lemma}
Идея в том, чтобы строго реализовать следующую схему: локально сжать вход в фиксированное малое окно и применить УАТ для MLP на конечномерном векторе.\\

\begin{proof}
\leavevmode
\underline{Шаг 0: конечномерность.}
Пространство $V_{\lambda,(T-1)\lambda L_{\mathrm{rf}}}$ конечномерно и отождествляется с $\mathbb{R}^N$,
где $N=d_V\,|Z_{(T-1)L_{\mathrm{rf}}}|$; аналогично $U_{\lambda,0}\cong\mathbb{R}^{d_U}$.
Следовательно, задача сводится к аппроксимации непрерывной функции $g\colon\mathbb{R}^N\to\mathbb{R}^{d_U}$
на компакте.

\underline{Шаг 1: <<сжатие>> (перенос всех координат в одно локальное место).}
Построим первые $T-2$ слоёв так, чтобы (на заданном компакте) они реализовывали отображение
\[
C\colon \ V_{\lambda,(T-1)\lambda L_{\mathrm{rf}}}\to W_{T-1}=V_{\lambda,L_{\mathrm{rf}}}(\text{с большим числом каналов}),
\]
которое кодирует все координаты входа в значениях признакового отображения вблизи нуля (внутри $Z_{L_{\mathrm{rf}}}$),
причём координаты кодируются \emph{по каналам}.
Формально: для каждой входной координаты $(\gamma,k)$ (где $\gamma\in Z_{(T-1)L_{\mathrm{rf}}}$)
выделим отдельный канал $c(\gamma,k)$ на последнем передвыходном слое и добьёмся, чтобы
\[
(C(\Phi))_{0,\,c(\gamma,k)}\approx \Phi_{\gamma,k},
\]
а остальные значения $C(\Phi)$ вне небольшого набора позиций можно считать несущественными.

Как это сделать свёрточными слоями фиксированной локальности?
Рассмотрим разложение $\gamma=\theta_1+\cdots+\theta_{T-1}$, где каждое $\theta_i\in Z_{L_{\mathrm{rf}}}$.
Такое разложение существует, потому что $\|\gamma\|_\infty\le (T-1)L_{\mathrm{rf}}$.
Тогда можно <<перетаскивать>> значение из позиции $\gamma$ к нулю за $T-1$ локальных шагов,
каждый раз используя сдвиг на $\theta_i$ внутри допустимого поля восприятия.

В линейном случае мы бы просто выбирали веса как индикаторы нужного сдвига и копировали координаты точно.
С нелинейностью $\sigma$ Яроцкий использует следующий стандартный приём:
координатные функционалы вида $\Phi\mapsto \Phi_{\eta,k}$ на конечномерном пространстве
аппроксимируются однослойными сетями, после чего линейные комбинации и <<маршрутизация>> по каналам
реализуются последующими свёрточными слоями. Это позволяет построить $C$ с любой требуемой точностью
на данном компакте.

\underline{Шаг 2: УА на локализованном векторе (последние два слоя).}
После <<сжатия>> $C(\Phi)$ зависит (с заданной точностью) от конечного набора скаляров, которые мы можем рассматривать
как вектор $x\in\mathbb{R}^N$, расположенный в каналах в одной точке решётки.
Последние два слоя при $\Lambda=0$ вычисляют выход в одной точке и имеют вид
полносвязной сети с одним скрытым слоем, построенной по этому вектору:
\[
(\widehat f_T\circ \widehat f_{T-1})(\Psi)
=
W^{(T)}\sigma(W^{(T-1)}\mathrm{vec}(\Psi)+b^{(T-1)})+b^{(T)}.
\]
По теореме универсальной аппроксимации (используемой Яроцким как теорема 2.1)
можно выбрать ширину скрытого слоя достаточно большой так, чтобы
эта композиция аппроксимировала произвольную непрерывную $g$ на компакте входных сигналов.

\underline{Шаг 3: композиция.}
Композиция <<первые слои $\approx C$>> и <<последние слои $\approx g\circ C^{-1}$ (на образе компакта)>>
даёт требуемую аппроксимацию $g$ базовой CNN с $\Lambda=0$.  
\end{proof}

Далее Яроцкий рассматривает сети с шагом субдескритизации $s$ и фиксированным $L_{\mathrm{rf}}$,
при этом накладывается условие достижимости (связности)
\begin{equation}
s\le 2L_{\mathrm{rf}}+1,
\label{eq:stride_condition}  
\end{equation}
которое гарантирует, что информация из каждой входной позиции может попасть к выходу через локальные окна \cite[(3.29)]{Yarotsky2018}.

\begin{definition}
Непрерывный функционал $f\colon V\to\mathbb{R}$ называется \emph{предельной точкой семейства свёрточных сетей с субдескритизацией},
если для любого компакта $K\subset V$, любого $\varepsilon>0$, любых $\lambda_0>0$, $\Lambda_0>0$
существует CNN с шагом субдескритизации $s$ (удовлетворяющим \eqref{eq:stride_condition}),
шагом сетки $\lambda\le\lambda_0$ и достаточно большой глубиной $T$, так что эффективный входной радиус
$\lambda L_{1,T}\ge \Lambda_0$ и
\[
\sup_{\Phi\in K}|f(\Phi)-\widehat f(\Phi)|<\varepsilon.
\]  
\end{definition}

\begin{theorem}
Функционал $f\colon V\to\mathbb{R}$ является предельной точкой семейства CNN с субдескритизацией
тогда и только тогда, когда $f$ непрерывен по норме.  
\end{theorem}

\begin{proof}
\textbf{(Необходимость).}
Точно как в теореме о характеризации предельных точек: берём сходящуюся последовательность $\Phi_n\to\Phi$,
рассматриваем компакт $K=\{\Phi\}\cup\{\Phi_n\}$ и аппроксимирующую $f$ сеть $\widehat f$ на $K$.
По непрерывности $\widehat f$ и оценке треугольником получаем непрерывность $f$.

\textbf{(Достаточность).}
Пусть $f$ непрерывен, $K\subset V$ компактен, $\varepsilon>0$.

\underline{Шаг 1: проектор на большое конечное окно.}
По лемме 1 можно выбрать $\lambda$ малым и параметры сети (в частности, глубину $T$)
так, чтобы проектор $P_{\lambda,\lambda L_{1,T}}$ был близок к тождеству на $K$:
\[
\sup_{\Phi\in K}\|P_{\lambda,\lambda L_{1,T}}\Phi-\Phi\|_V<\delta,
\]
а равномерная непрерывность $f$ на компакте $K$ даёт
\begin{equation}
\sup_{\Phi\in K}|f(P_{\lambda,\lambda L_{1,T}}\Phi)-f(\Phi)|<\varepsilon/2.
\label{eq:disc_cutoff}  
\end{equation}

\underline{Шаг 2: конечномерная аппроксимация.}
Отображение $g\coloneq f\circ P_{\lambda,\lambda L_{1,T}}$ действует на конечномерном пространстве
$V_{\lambda,\lambda L_{1,T}}\cong\mathbb{R}^N$, и достаточно аппроксимировать $g$ на компакте
$P_{\lambda,\lambda L_{1,T}}K$.

\underline{Шаг 3: реализуемость $g$ архитектурой с субдескритизацией.}
Благодаря условию \eqref{eq:stride_condition} любая входная позиция на мелкой решётке
может быть выражена как $s\gamma+\theta$ с $\theta\in Z_{L_{\mathrm{rf}}}$,
а значит локальный слой с шагом $s$ может <<забирать>> информацию из каждой позиции и переносить её на более грубую решётку.
Итерируя по слоям, можно построить (с достаточной шириной каналов) схему <<сжатия>>,
аналогичную доказательству Леммы \ref{lem:approx_degen_window}, которая кодирует конечномерный вход
в признаки в небольшой области на последнем слое.
После этого последние слои реализуют MLP с одним скрытым слоем над полученным конечномерным вектором признаков,
и по УАТ аппроксимируют произвольное непрерывное $g$ на компакте.

То есть существует downsampling-CNN $\widehat f$ такая, что
\begin{equation}
\sup_{\Phi\in K}\big|\widehat f(\Phi)-f(P_{\lambda,\lambda L_{1,T}}\Phi)\big|<\varepsilon/2.
\label{eq:finite_approx}  
\end{equation}
Совмещая \eqref{eq:disc_cutoff} и \eqref{eq:finite_approx}, получаем
$\sup_{\Phi\in K}|\widehat f(\Phi)-f(\Phi)|<\varepsilon$.  
\end{proof}

Эти две теоремы дают строгую характеризацию аппроксимационной мощности CNN
в двух режимах:
\begin{itemize}
  \item \textbf{без пулинга:} $\overline{\mathcal{H}}_{\mathrm{CNN}} = \{\text{непрерывные и сдвиг-эквивариантные } f\colon V\to U\};$
  \item \textbf{с субдескритизацией:} $\overline{\mathcal{H}}_{\mathrm{CNN}} = \{\text{непрерывные функционалы } f\colon V\to\mathbb{R}\}.$
\end{itemize}
В первом случае симметрия действует как \emph{жёсткое ограничение} на возможный предел,
во втором --- субдескритизация разрушает эквивариантность в пределе, и остаётся только топологическое условие непрерывности.

\subsubsection{Стандартная модель Чжоу: пространство гипотез и универсальность}
В статье Чжоу \cite{Zhou2018} рассматривается \emph{чисто свёрточная} глубокая сеть без полносвязных матриц:
все линейные преобразования имеют тёплицеву (Toeplitz) структуру, соответствующую одномерной свёртке
(вектор $x\in\mathbb R^d$ интерпретируется как 1D-сигнал длины $d$).

\begin{definition}[Маска фильтра и матрица свёртки]
Фиксируем длину фильтра $s\ge 2$.
\emph{Маска фильтра} (filter mask) --- это последовательность $w=(w_k)_{k\in\mathbb Z}$, причём
\[
w_k=0\quad \text{для всех } k\notin\{0,1,\dots,s\}.
\]
Для вектора $v=(v_1,\dots,v_D)^\top\in\mathbb R^D$ (с нулевым продолжением вне $\{1,\dots,D\}$) определим свёртку
\[
(w*v)_i=\sum_{k=1}^{D} w_{i-k}\,v_k .
\]
Этому соответствует \emph{матрица свёртки} $T=T(w)$ тёплицева типа размера $(D+s)\times D$ с постоянными диагоналями,
то есть $Tv=w*v$.  
\end{definition}

\begin{definition}[Глубокая CNN \cite{Zhou2018}]
Пусть глубина равна $J$. Вводим ширины
\[
d_0=d,\qquad d_j=d+js,\quad j=1,\dots,J.
\]
Пусть $T^{(j)}=T(w^{(j)})$ --- тёплицева матрица свёртки размера $d_j\times d_{j-1}$, а $\sigma(u)=u_+=\max\{u,0\}$
--- ReLU. Определим промежуточные представления (feature maps)
\[
h^{(0)}(x)=x,\qquad
h^{(j)}(x)=\sigma\big(T^{(j)}h^{(j-1)}(x)-b^{(j)}\big),\quad j=1,\dots,J,
\]
где $b^{(j)}\in\mathbb R^{d_j}$ --- векторы смещений.
\end{definition}

\begin{definition}[Ограничение на смещения (bias pattern)]
Для $j\le J-1$ в статье берётся специальная форма:
компоненты $b^{(j)}_{s+1},\dots,b^{(j)}_{d_j-s}$ равны одному и тому же числу (середина повторяется),
а крайние $2s$ компонент могут быть произвольными.
(Именно это позволяет <<выключать>>/<<включать>> нужные координаты в конструкциях.)
\end{definition}

\begin{definition}[Гипотезное пространство]
Выходом сети служит линейная комбинация компонент последнего слоя:
\[
\mathcal H^{J}_{w,b}
=
\left\{
x\mapsto \sum_{k=1}^{d_J} c_k\,h^{(J)}_k(x)\ :\ c\in\mathbb R^{d_J}
\right\}.
\]
\end{definition}
Это и есть класс функций, реализуемых \emph{глубокой свёрточной сетью Чжоу} глубины $J$ в \cite{Zhou2018}.

Два ключевых факта обеспечивающих конструктивную выразительность класса $\mathcal H^{J}_{w,b}$:
\begin{enumerate}
  \item произведение тёплицевых матриц снова является тёплицевой матрицей, отвечающей свёртке свёрток;
  \item любую конечную маску можно разложить в свёрточное произведение коротких масок длины $s+1$.
\end{enumerate}

\begin{lemma}
\label{lem:toeplitz_comp}
Пусть $w^{(1)},\dots,w^{(J)}$ --- векторы с компонентами из $\{0,\dots,s\}$ и $T^{(j)}=T(w^{(j)})$.
Тогда существует (конечная) последовательность $W$ такая, что
\[
T^{(J)}T^{(J-1)}\cdots T^{(1)}=T(W),
\]
и при этом $W=w^{(J)}*w^{(J-1)}*\cdots*w^{(1)}$ (ассоциативно).  
\end{lemma}
\begin{proof}
Для любого вектора $v$ имеем $T(w)v=w*v$ по определению.
Тогда
\[
T^{(2)}T^{(1)}v = T(w^{(2)})(w^{(1)}*v)=w^{(2)}*(w^{(1)}*v)=(w^{(2)}*w^{(1)})*v.
\]
Индукцией по $J$ получаем
\(
T^{(J)}\cdots T^{(1)}v=(w^{(J)}*\cdots*w^{(1)})*v
\)
для всех $v$, следовательно матрицы совпадают.  
\end{proof}

\begin{theorem}[Cвёрточная факторизация конечной маски, теорема C \cite{Zhou2018}]
\label{thm:zhou_C}
Пусть $s\ge 2$ и последовательность $W=(W_k)_{k\in\mathbb Z}$ из векторов с компонентами из $\{0,1,\dots,M\}$.
Тогда существуют $J<\frac{M}{s-1}+1$ и маски $w^{(1)},\dots,w^{(J)}$,
каждая состоящая из элементов $\{0,\dots,s\}$, такие что
\[
W=w^{(J)}*w^{(J-1)}*\cdots*w^{(1)}.
\]  
\end{theorem}
\begin{proof}
Следуем доказательству из оригинальной статьи.
Для конечной неотрицательной последовательности $a$ введём \emph{символ}
\[
\widehat a(z)=\sum_{k\ge 0} a_k z^k,
\]
это многочлен степени $\deg \widehat a=\max\{k:a_k\ne 0\}$.
Для свёртки выполняется тождество символов:
\[
\widehat{a*b}(z)=\widehat a(z)\,\widehat b(z),
\]
поскольку коэффициенты произведения многочленов равны свёртке коэффициентов.

Пусть $\widehat W(z)$ --- многочлен степени $M$ с вещественными коэффициентами.
Разложим его на линейные и квадратичные множители над $\mathbb R$.
Комплексные корни идут попарно: если $z_k=x_k+iy_k$ корень, то и $\overline z_k=x_k-iy_k$ корень, и
\[
(z-z_k)(z-\overline z_k)=z^2-2x_k z + (x_k^2+y_k^2),
\]
то есть вещественный квадратичный множитель. Поэтому существует факторизация вида
\[
\widehat W(z)=W_M
\prod_{k=1}^{K}\big(z^2-2x_k z+(x_k^2+y_k^2)\big)\,
\prod_{k=2K+1}^{M}(z-x_k),
\]
где второе произведение отвечает вещественным корням (с учётом кратностями).

Теперь сгруппируем множители в блоки так, чтобы степень каждого блока была $\le s$:
\begin{itemize}
\item берём до $\lfloor s/2\rfloor$ квадратичных множителей (или $\lfloor (s-1)/2\rfloor$ квадратичных + 1 линейный),
\item либо берём до $s$ линейных множителей,
\end{itemize}
пока не исчерпаем весь набор. Тогда получим разложение
\[
\widehat W(z)=\widehat w^{(J)}(z)\cdots \widehat w^{(2)}(z)\widehat w^{(1)}(z),
\quad \deg \widehat w^{(j)}\le s,
\]
где каждый $\widehat w^{(j)}$ имеет вещественные коэффициенты.
Переходя от многочленов к последовательностям коэффициентов, получаем
$W=w^{(J)}*\cdots*w^{(1)}$,
а оценка на $J$ следует из того, что каждый блок уменьшает суммарную степень минимум на $(s-1)$, кроме последнего блка.  
\end{proof}

\begin{theorem}[Универсальность глубокой CNN, теорема A\cite{Zhou2018}]
Пусть $2\le s\le d$ и $\Omega\subset\mathbb R^d$ --- компакт.
Для любого $f\in C(\Omega)$ существуют глубины $J\to\infty$, наборы масок $w$, смещений $b$ и функции
$g_J\in \mathcal H^{J}_{w,b}$ такие, что
\[
\|f-g_J\|_{C(\Omega)}\to 0\quad \text{при } J\to\infty.
\]  
\end{theorem}
\begin{proof}
Доказательство оформим конструктивно, следуя плану Чжоу и стандартной аппроксимационной теории.

\underline{Шаг 1: редукция к полиномам.}
По теореме Стоуна--Вейерштрасса алгебра полиномов на $\mathbb R^d$ плотна в $C(\Omega)$.
Значит, для любого $\varepsilon>0$ существует полином $p$ такой, что
\[
\|f-p\|_{C(\Omega)}<\varepsilon/3.
\]

\underline{Шаг 2: разложение полинома на ridge-термы.}
Нам достаточно доказать, что любой многочлен $p$ можно представить как конечную сумму
\[
p(x)=\sum_{q=1}^{Q} p_q(\langle \xi_q,x\rangle),
\qquad \xi_q\in\mathbb R^d,
\]
где $p_q$ --- одномерные полиномы.
Это следует из того, что пространство полиномов степени $\le n$ конечномерно,
а семейство функций $\{(\langle \xi,x\rangle)^k:\ \xi\in\mathbb R^d,\ 0\le k\le n\}$
порождает его как линейную оболочку: действительно, для фиксированного $n$ рассмотрим конечный набор направлений
$\xi_1,\dots,\xi_R$, выбранный так, чтобы матрица значений мономов на наборе направлений имела полный ранг
(например, берём достаточно много <<общих>> $\xi_r$; тогда соответствующая вандермондова система невырождена).
Тогда любой моном $x^\alpha$ степени $|\alpha|\le n$ выражается как линейная комбинация степеней
$(\langle \xi_r,x\rangle)^{|\alpha|}$, а значит и любой полином $p$ выражается как сумма одномерных полиномов от $\langle\xi_r,x\rangle$.

\underline{Шаг 3: реализация ridge-термов сетью Чжоу.}
Зафиксируем один терм $p_q(\langle\xi_q,x\rangle)$. Сначала объясним, как получить линейную форму $\langle\xi_q,x\rangle$
внутри сети, затем как реализовать одномерный полином/его аппроксимацию через ReLU.

\emph{(3a) Линейная форма как одна координата $T(W)x$.}
Пусть нам нужно получить набор скаляров $\alpha_0\cdot x,\alpha_1\cdot x,\dots,\alpha_m\cdot x$
(в частности, можно взять $\alpha_0=\xi_q$ и дополнительные $\alpha_k$ для построения ломаной аппроксимации).
Сформируем конечную последовательность $W$ длины $(m+1)d$ <<склейкой>> коэффициентов
так, чтобы некоторые строки матрицы $T(W)$ совпали с $\alpha_k^\top$.
По теореме \ref{thm:zhou_C} найдётся факторизация $W=w^{(J)}*\cdots*w^{(1)}$ с масками длины $s+1$ и глубиной $J=O(m)$.
По лемме \ref{lem:toeplitz_comp} получаем
\[
T^{(J)}\cdots T^{(1)}=T(W),
\]
то есть произведение тёплиц-слоёв вычисляет нужные линейные формы как компоненты $T(W)x$.

\emph{(3b) ReLU-шарниры и линейные сплайны.}
Классический факт: любую непрерывную $u\colon[a,b]\to\mathbb R$ можно аппроксимировать линейным сплайном,
а линейный сплайн представим через <<шарниры>> $(t-\tau)_+=\sigma(t-\tau)$:
если узлы $\tau_1<\dots<\tau_m$, то
\[
S(t)=\beta_0+\beta_1 t+\sum_{k=1}^{m} \gamma_k (t-\tau_k)_+,
\tag{$\star$}\label{eq:hinge_spline}
\]
где коэффициенты $\gamma_k$ выражаются через конечные разности наклонов на соседних интервалах.
Это так, потому что производная $S'(t)$ кусочно-постоянна и испытывает скачки $\gamma_k$ в $\tau_k$,
а интегрирование даёт \eqref{eq:hinge_spline} (смотри например \cite[Глава 5]{DeVore1993}).

Следовательно, чтобы реализовать/аппроксимировать $p_q(t)$ на множестве значений $t=\langle\xi_q,x\rangle$,
достаточно реализовать множество функций $(t-\tau_k)_+$ и затем взять линейную комбинацию.

\emph{(3c) Встраивание шарниров в $\mathcal H^J_{w,b}$.}
После шага (3a) нужный скаляр $t=\langle\xi_q,x\rangle$ появляется как одна из координат некоторого $h^{(J-1)}(x)$
(либо прямо как координата $T(W)x$ до применения ReLU).
Затем на последнем слое мы выбираем смещения $b^{(J)}$ так, чтобы отдельные координаты имели вид
\[
h^{(J)}_{\ell}(x)=\sigma\big(t-\tau_k\big)=(t-\tau_k)_+
\]
для нужных $\tau_k$.
После этого линейная комбинация с коэффициентами $c_\ell$ реализует \eqref{eq:hinge_spline}.

Итак, для каждого ridge-терма можно построить $g_q\in \mathcal H^J_{w,b}$, аппроксимирующий
$p_q(\langle\xi_q,x\rangle)$ с ошибкой $<\varepsilon/(3Q)$ на $\Omega$.
Суммируя по $q$ (линейность выхода по $c$), получаем $g=\sum_{q=1}^{Q} g_q\in \mathcal H^J_{w,b}$ и
\[
\|p-g\|_{C(\Omega)}<\varepsilon/3.
\]
Вместе с шагом 1:
\[
\|f-g\|_{C(\Omega)}\le \|f-p\|_{C(\Omega)}+\|p-g\|_{C(\Omega)}<2\varepsilon/3.
\]
Уточняя точности на шагах 1 и 3 (заменой $\varepsilon/3$ на $\varepsilon/2$ и тому подобное), получаем требуемое $\|f-g\|_{C(\Omega)}<\varepsilon$.
\end{proof}

Далее фиксируем $d=1$, $\Omega=[0,1]$, нас будет интересовать количественная оценка аппроксимации липшицевых функций.
\begin{theorem}
Пусть $d=1$, $f\in C([0,1])$ липшицева с полу-нормой
\(
\|f\|_{\mathrm{Lip}}=\sup_{x\ne y}\frac{|f(x)-f(y)|}{|x-y|}.
\)
Тогда существует константа $c>0$, не зависящая от $f$ и $J$, такая что
\[
\inf_{g\in \mathcal H^{J}_{w,b}}\|f-g\|_{\infty}\le c\,\|f\|_{\mathrm{Lip}}\,J^{-1/2}.
\]  
\end{theorem}
\begin{proof}
\emph{Шаг 1: линейный сплайн на равномерной сетке.}
Для $m\in\mathbb N$ рассмотрим узлы $x_k=k/m$, $k=0,\dots,m$, и линейный интерполянт $S_m f$ по этим узлам.
Для любого $x\in[x_k,x_{k+1}]$ по липшицевости
\[
|f(x)-S_m f(x)|\le \frac{\|f\|_{\mathrm{Lip}}}{2m},
\]
откуда
\begin{equation}
\|f-S_m f\|_\infty\le \frac{\|f\|_{\mathrm{Lip}}}{2m}.
\label{eq:spline_err}  
\end{equation}

\emph{Шаг 2: представление сплайна через шарниры ReLU.}
Из общей формулы \eqref{eq:hinge_spline} следует: существует представление
\begin{equation}
S_m f(x)=\beta_0+\beta_1 x+\sum_{k=1}^{m-1}\gamma_k (x-x_k)_+,
\label{eq:spline_hinge_1d}  
\end{equation}
где $x_k=k/m$, а коэффициенты $\gamma_k$ выражаются через скачки наклонов $S_m f$ в узлах.

\emph{Шаг 3: реализация \eqref{eq:spline_hinge_1d} в классе Чжоу при глубине $J\gtrsim m^2$.}
Поясним конструктивно, почему при фиксированной длине фильтра $s+1$ можно реализовать
семейство шарниров $(x-x_k)_+$ для $k=1,\dots,m-1$ внутри $\mathcal H^J_{w,b}$, если $J$ достаточно велик.

В 1D вход $x$ является скаляром. С помощью тёплиц-слоёв можно <<размножить>> информацию об $x$ по разным координатам
и затем, используя разные смещения в последнем слое, получить набор функций
\[
x\longmapsto (x-x_k)_+,\qquad k=1,\dots,m-1.
\]
Технически это реализуется так же, как в явной конструкции для получения координат
$(\alpha_k\cdot x-t_k)_+$: линейная стадия создаёт нужные аффинные формы, а последний ReLU + смещение превращает их в шарниры.
Ограничение на смещение не мешает, поскольку <<особые>> компоненты допускаются на краях,
а середину можно держать константной (используя повторяющийся параметр).

Почему возникает масштаб $J\gtrsim m^2$?
Интуитивно: чтобы получить $m$ разных <<переломов>> на фиксированном фильтре, нужно достаточно много итераций локального
распространения/кодирования.
Следовательно, при $J\ge C m^2$ существует $g\in\mathcal H^J_{w,b}$ такой, что $g(x)=S_m f(x)$ (или
$\|g-S_m f\|_\infty$ поглощается константой).

\emph{Шаг 4: выбор $m\asymp \sqrt{J}$ и итог.}
Положим $m=\lfloor c_0\sqrt{J}\rfloor$ так, чтобы $J\ge C m^2$.
Тогда из \eqref{eq:spline_err} и равенства/приближения $g\approx S_m f$ получаем
\[
\|f-g\|_\infty
\le \|f-S_m f\|_\infty + \|S_m f-g\|_\infty
\le \frac{\|f\|_{\mathrm{Lip}}}{2m}+o(1)
\le c\,\|f\|_{\mathrm{Lip}}\,J^{-1/2}.
\]  
\end{proof}

\begin{remark}
В оригинальной статье основной количественный результат формулируется для более гладких классов (Соболевские пространства)
и даёт более быстрые скорости на 1D при дополнительной регулярности.
Теорема служит иллюстрацией подхода \emph{кусочно-линейная аппроксимация + ReLU-шарниры + тёплиц-структура}.  
\end{remark}

\subsubsection{Подход Ли-Лин-Шеня: универсальность полносвязных свёрточных сетей для сдвиг-инвариантных функций}

\textbf{Контекст.}
В предыдущих фрагментах мы видели, что \emph{симметрийные ограничения} (например, сдвиг-эквивариантность) радикально меняют ответ на вопрос об универсальной аппроксимации.
В этой части рассматривается \emph{сдвиг-инвариантная} постановка на конечной периодической решётке (дискретный тор) и даётся \emph{точный} критерий, когда полноcвязные свёрточные сети (FCNN) с \emph{фиксированным} числом каналов и \emph{маленьким} ядром (размер 2) обладают свойством УА.

Зафиксируем размерность $d\in\mathbb{N}$ и вектор размеров $n=(n_1,\dots,n_d)$.
Рассмотрим пространство сигналов
\[
X \coloneq  \mathbb{R}^{n_1\times\cdots\times n_d}\cong \mathbb{R}^{|n|},\qquad |n|\coloneq \prod_{s=1}^d n_s.
\]
Индексы будем понимать по модулю $n_s$ (периодические граничные условия).

Для $k=(k_1,\dots,k_d)$, где $k_s\in\{0,1,\dots,n_s-1\}$, определим \emph{оператор сдвига} $T_k\colon X\to X$:
\[
[T_k x]_i \coloneq  x_{i-k}\quad(\text{вычитание по модулю }n).
\]
Действие $k\mapsto T_k$ образует (конечную) группу сдвигов.

В данном контексте уточним раннее введённые понятия.
\begin{itemize}
  \item Функция $F\colon X\to\mathbb{R}$ называется \emph{сдвиг-инвариантной}, если
  \[
    F(T_k x)=F(x)\quad\forall x\in X,\ \forall k.
  \]
  \item Отображение $\Phi\colon X\to X$ называется \emph{сдвиг-эквивариантным}, если
  \[
    \Phi(T_k x)=T_k\Phi(x)\quad\forall x\in X,\ \forall k.
  \]
\end{itemize}

Для $w,x\in X$ (ядро $w$ и сигнал $x$) \emph{свёртка} $w*x\in X$ определяется стандартно (индексы периодические):
\[
[w*x]_i\coloneq \sum_{j} w_j\,x_{i-j}.
\]
Из определения следует коммутирование со сдвигом:
\[
T_k(w*x)=(T_k w)*x=w*(T_k x).
\]
Это и есть формальная причина, почему свёрточные слои при \emph{разделении весов} естественно порождают сдвиг-эквивариантные/инвариантные модели.

Обсудим подробнее архитектуру FCCN.
Далее $\sigma\colon \mathbb{R}\to\mathbb{R}$ --- нелинейность (в части про неостаточные сети будет важна $\sigma=\mathrm{ReLU}$).

\begin{definition}[Носитель тензора \cite{Li_Lin_Shen2022}]
Определим \emph{носитель} тензора $x\in X$ как минимальный $d$-вектор $\mathrm{supp}(x)=(j_1,\dots,j_d)$, такой что
если для мультииндекса $i=(i_1,\dots,i_d)$ выполнено $i_s>j_s$ хотя бы для одного $s$, то $[x]_i=0$.
Пишем $\mathrm{supp}(x)\le \ell$ (покомпонентно), если $j_s\le \ell_s$ для всех $s$.  
\end{definition}

\textit{Интерпретация}: ядро с $\ell=(2,\dots,2)$ использует ближайших соседей <<в одном шаге>> по каждой координате.

Пусть $r\in\mathbb{N}$ --- число каналов признаков. Для фиксированного ограничения $\ell$ на носитель ядра зададим класс функций
\begin{equation}
\mathcal{F}_{r,\ell}
\coloneq \left\{
x\mapsto \sum_{i=1}^r v_i\,\sigma\!\bigl(w_i*x+b_i\mathbf{1}\bigr)
:\ w_i\in X,\ \mathrm{supp}(w_i)\le \ell,\ v_i,b_i\in\mathbb{R}
\right\}.
\label{eq:Frl}
\end{equation}
Здесь $\mathbf{1}\in X$ --- тензор из единиц (смещение одинаково по всем координатам), а коэффициенты $v_i$ <<смешивают>> каналы.

Пусть $g\colon X\to\mathbb{R}$ --- \emph{липшицева} и \emph{инвариантная относительно перестановок координат} функция:
значение $g(x)$ не меняется при любой перестановке компонент $x$.
Например, сумма $g(x)=\sum_i [x]_i$ или максимум $g(x)=\max_i [x]_i$ (\emph{глобальный пулинг}).


Определим множества
\begin{align}
\mathrm{CNN}_{r,\ell}
&\coloneq \left\{ g\circ f_m\circ\cdots\circ f_1:\ f_1,\dots,f_m\in\mathcal{F}_{r,\ell},\ m\ge 1\right\},\label{eq:CNN}\\
\mathrm{resCNN}_{r,\ell}
&\coloneq \left\{ g\circ (\mathrm{id}+f_m)\circ\cdots\circ (\mathrm{id}+f_1):\ f_1,\dots,f_m\in\mathcal{F}_{r,\ell},\ m\ge 1\right\}.\label{eq:resCNN}
\end{align}
По построению любые элементы $\mathrm{CNN}_{r,\ell}$ и $\mathrm{resCNN}_{r,\ell}$ являются сдвиг-инвариантными функциями $X\to\mathbb{R}$ (периодичность и разделение весов сохраняют симметрию на каждом слое).

\begin{definition}
Функциональное семейство $\mathcal{H}\subset\{X\to\mathbb{R}\}$ обладает \emph{сдвиг-инвариантной универсальной аппроксимацией в $L^p$},
если:
\begin{enumerate}
\item каждая $H\in\mathcal{H}$ сдвиг-инвариантна;
\item для любой сдвиг-инвариантной непрерывной (или $L^p$-) функции $F\colon X\to\mathbb{R}$, любого компакта $K\subset X$, любого $\varepsilon>0$ и $p\in[1,\infty)$
существует $H\in\mathcal{H}$, такое что $\|F-H\|_{L^p(K)}\le \varepsilon$.
\end{enumerate}  
\end{definition}

\begin{theorem}
\label{thm:lls_main}
Имеют место утверждения \cite{Li_Lin_Shen2022}:
\begin{enumerate}
\item (\emph{Достаточность:}) Остаточное гипотезное пространство $\mathrm{resCNN}_{r,\ell}$ обладает сдвиг-инвариантной УА при $r\ge 1$ и $\ell\ge 2$ (покомпонентно).
Пространство $\mathrm{CNN}_{r,\ell}$ обладает сдвиг-инвариантной УА при $r\ge 2$ и $\ell\ge 2$.
\item (\emph{Оптимальность размера ядра:}) Если $\min_s \ell_s=1$, то ни $\mathrm{resCNN}_{\infty,\ell}+\mathbb{R}$\footnote{Примечание: как и в алгебре это множество всевозможных сумм}, ни $\mathrm{CNN}_{\infty,\ell}+\mathbb{R}$ не обладают сдвиг-инвариантной УА.
\item (\emph{Оптимальность числа каналов:}) Семейство $\mathrm{CNN}_{1,\infty}+\mathbb{R}$ не обладает сдвиг-инвариантной УА.
\end{enumerate}  
\end{theorem}
\begin{proof}
\underline{Шаг 1: из сдвиг-эквивариантной получить сдвиг-инвариантную УА.}
Сначала временно \emph{уберём} финальный пулинг и рассмотрим отображения $X\to X$ (пиксельный выход), то есть эквивариантную версию.
Определим
\[
\mathrm{eqCNN}_{r,\ell}\coloneq \{f_m\circ\cdots\circ f_1:\ f_i\in\mathcal{F}_{r,\ell}\},\qquad
\mathrm{eqresCNN}_{r,\ell}\coloneq \{(\mathrm{id}+f_m)\circ\cdots\circ(\mathrm{id}+f_1)\}
\]
и соответствующее понятие \emph{сдвиг-эквивариантной УА} (аппроксимация эквивариантных отображений $X\to X$ в $L^p(K)$).

\begin{proposition}[Связь инвариантной и эквивариантной УА через пулинг]
\label{prop:pooling_bridge}
Пусть $g\colon X\to\mathbb{R}$ --- липшицева, инвариантна относительно перестановок координат и $g(X)=\mathbb{R}$.
Если семейство отображений $\mathcal{A}\subset\{X\to X\}$ обладает сдвиг-эквивариантной УА, то семейство функций
\[
\mathcal{B}\coloneq \{g\circ \varphi:\ \varphi\in\mathcal{A}\}\subset\{X\to\mathbb{R}\}
\]
обладает сдвиг-инвариантной УА.
\end{proposition}
\begin{proof}[Доказательство утверждения]
Пусть $F\colon X\to\mathbb{R}$ --- сдвиг-инвариантная $L^p$-функция, $K\subset X$ --- компакт.
Можно считать (расширив компакт при необходимости), что $K=[-a,a]^n$.

Определим подмножество
\[
K_1\coloneq \{x\in K:\ [x]_1>[x]_i\ \forall i\neq 1\}.
\]
Тогда (с точностью до множества меры нуль) $K=\bigcup_k T_k K_1$ и эти множества дизъюнктны.

Положим $\varepsilon'\coloneq \varepsilon/(|n|(1+\mathrm{Lip}(g)))$.
По универсальности для \emph{несимметричного} случая (\cite[Теорема 3.8]{Li_Lin_Shen2022} используется в статье как внешний инструмент),
существует (не обязательно эквивариантное) $u\colon X\to X$ такое, что
\begin{equation}
\|F-g\circ u\|_{L^p(K)}\le \varepsilon'.
\label{eq:bridge_27}
\end{equation}
Далее сделаем $u$ эквивариантной.

Так как $u\in L^p$, найдётся компакт $O\subset K_1$ с $\|u\|_{L^p(K_1\setminus O)}\le \varepsilon'$.
Возьмём гладкую финитную функцию $\chi$ со значениями в $[0,1]$, такую что $\chi|_O=1$ и $\chi|_{K_1^c}=0$.
Положим $\tilde u\coloneq \chi u$.

Для $x\in T_kK_1$ определим
\[
f(x)\coloneq T_k\bigl(\tilde u(T_{-k}x)\bigr),
\]
а на дополнении $\bigcup_k T_kK_1$ положим $f(x)\coloneq 0$.
Поскольку $\chi$ зануляется на границе $K_1$, $f$ непрерывно, и прямой проверкой получаем её сдвиг-эквивариантность.

Оценим ошибку. Так как $F$ и $g\circ f$ сдвиг-инвариантны, достаточно оценивать на $K_1$:
\begin{equation}
\|F-g\circ f\|_{L^p(K)} = |n|\,\|F-g\circ f\|_{L^p(K_1)}.
\label{eq:bridge_28}
\end{equation}
На $K_1$ имеем $f|_{K_1}=\tilde u$, а значит
\[
\|u-f\|_{L^p(K_1)}=\|u-\tilde u\|_{L^p(K_1)}\le \|u\|_{L^p(K_1\setminus O)}\le \varepsilon'.
\]
Липшицевость $g$ даёт $\|g\circ u-g\circ f\|_{L^p(K_1)}\le \mathrm{Lip}(g)\varepsilon'$,
и вместе с \eqref{eq:bridge_27}--\eqref{eq:bridge_28} получаем
\[
\|F-g\circ f\|_{L^p(K)}\le (1+\mathrm{Lip}(g))|n|\varepsilon'=\varepsilon.
\]
Осталось применить эквивариантную УА для семейства $\mathcal{A}$ к эквивариантному $f$: существует $\varphi\in\mathcal{A}$,
$\|f-\varphi\|_{L^p(K)}\le \varepsilon/\mathrm{Lip}(g)$, и тогда $\|F-g\circ\varphi\|_{L^p(K)}\le \varepsilon$.  
\end{proof}
\underline{Шаг 2: динамическая интерпретация ResNet и аппроксимация потоков остаточными блоками.}
Ключевая идея: композицию остатков
\(
(\mathrm{id}+f_m)\circ\cdots\circ(\mathrm{id}+f_1)
\)
удобно понимать как (грубую) дискретизацию потока ОДУ, а значит удобно ввести непрерывно-временной <<идеальный>> класс аппроксиматоров.

Пусть $f\colon X\to X$ липшицево. Для $T\in\mathbb{R}$ определим \emph{поток} $\phi(f,T)\colon X\to X$ как
\[
\phi(f,T)(x)=z(T),\qquad \dot z(t)=f(z(t)),\ z(0)=x.
\]
Стандартные результаты теории ОДУ дают, что $\phi(f,T)$ липшицево (и даже билипшицево, с обратным $\phi(-f,T)$).
Определим
\begin{align}
\mathrm{eqCODE}_{r,\ell}
&\coloneq \left\{\phi(f_m,t_m)\circ\cdots\circ\phi(f_1,t_1):\ f_i\in\mathcal{F}_{r,\ell},\ t_i\in\mathbb{R},\ m\ge 1\right\},\\
\mathrm{CODE}_{r,\ell}
&\coloneq \left\{g\circ \Phi:\ \Phi\in \mathrm{eqCODE}_{r,\ell}\right\}.
\end{align}

\begin{proposition}
\label{prop:euler_resnet}
Пусть $f:X\to X$ липшицево с константой $L$ и $K\subset X$ --- компакт.
Тогда для любого $T\in\mathbb{R}$ и любого $\varepsilon>0$ существует $N\in\mathbb{N}$ такое, что
\[
\sup_{x\in K}\left\|\phi(f,T)(x)-\left(\mathrm{id}+\frac{T}{N}f\right)^{N}(x)\right\|\le \varepsilon,
\]
где $\circ N$ --- $N$-кратная композиция.
В частности, если $f\in\mathcal{F}_{r,\ell}$, то поток $\phi(f,T)$ равномерно на $K$ аппроксимируется элементами $\mathrm{eqresCNN}_{r,\ell}$.  
\end{proposition}
\begin{proof}[Доказательство утверждения]
Это стандартная оценка сходимости явной схемы Эйлера для ОДУ на конечном горизонте времени.
Пусть $h\coloneq T/N$ (для $T<0$ рассуждение то же, так как можно заменить $f\mapsto -f$ и $T\mapsto -T$).
Рассмотрим истинную траекторию $z(t)$ и дискретную $x_k$:
\[
x_{k+1}=x_k+h f(x_k),\qquad x_0=x,\qquad t_k\coloneq kh.
\]
По формуле Тейлора с интегральным остатком для ОДУ получаем локальную ошибку порядка $h^2$:
\[
z(t_{k+1}) = z(t_k)+h f(z(t_k)) + \int_{t_k}^{t_{k+1}}(f(z(s))-f(z(t_k)))\,ds.
\]
Обозначим $e_k\coloneq z(t_k)-x_k$.
Тогда
\[
e_{k+1}=e_k+h(f(z(t_k))-f(x_k)) + r_k,
\qquad r_k\coloneq \int_{t_k}^{t_{k+1}}(f(z(s))-f(z(t_k)))\,ds.
\]
Липшицевость даёт $\|f(z(t_k))-f(x_k)\|\le L\|e_k\|$.
Далее, так как $z$ липшицева по времени на $[0,T]$ (оценка через $\sup\|f\|$ на компактной трубке вокруг траекторий),
получаем $\|r_k\|\le C h^2$ для некоторой константы $C$ (зависящей от $K,f,T$). Тогда
\[
\|e_{k+1}\|\le (1+Lh)\|e_k\| + Ch^2.
\]
Итерацией и дискретным неравенством Гронуолла:
\[
\|e_N\|\le Ch^2\sum_{j=0}^{N-1}(1+Lh)^j
= Ch^2\,\frac{(1+Lh)^N-1}{(1+Lh)-1}
\le \frac{C}{L}\,h\,(e^{LT}-1).
\]
Правая часть $\to 0$ при $h\to 0$, то есть при $N\to\infty$.
Причём все оценки можно сделать равномерными по $x\in K$ (за счёт компактности $K$ и непрерывности $f$).
Отсюда следует требуемое равномерное приближение потока композицией $(\mathrm{id}+h f)$.

Если $f\in\mathcal{F}_{r,\ell}$, то отображение $x\mapsto x+h f(x)$ --- это ровно один остаточный блок
$(\mathrm{id}+f_h)$ с $f_h\coloneq h f\in\mathcal{F}_{r,\ell}$, а значит $N$-кратная композиция лежит в $\mathrm{eqresCNN}_{r,\ell}$.
\end{proof}

Дальнейший решающий факт заключается в том, что при минимальных параметрах $(r,\ell)=(1,2)$ динамический класс уже универсален:
\begin{theorem}
\label{thm:eqcode_UA}
Динамическое гипотезное пространство $\mathrm{eqCODE}_{1,2}$ обладает сдвиг-эквивариантной УА, а значит $\mathrm{CODE}_{1,2}$ обладает сдвиг-инвариантной УА.  
\end{theorem}
Её доказательство опирается на <<перестановку области>> (domain rearrangement)
через композиции потоков с маленьким ядром и на схему <<coordinate zooming + point matching>>.

\underline{Шаг 3: завершение доказательства пункта (1) теоремы.}\\
\emph{(а) Остаточный случай.}
По Теореме \ref{thm:eqcode_UA} класс $\mathrm{eqCODE}_{1,2}$ эквивариантно-универсален.
По предложению \ref{prop:euler_resnet} каждый поток $\phi(f,t)$ аппроксимируется композициями остаточных блоков, то есть
$\mathrm{eqresCNN}_{1,2}$ наследует сдвиг-эквивариантную УА (аппроксимация каждой $\phi(f_i,t_i)$ и устойчивость относительно композиции на компактах).
Далее по предложению \ref{prop:pooling_bridge} получаем сдвиг-инвариантную УА для $\mathrm{resCNN}_{1,2}$.
Для общих $r\ge 1$, $\ell\ge 2$ имеем включение $\mathrm{resCNN}_{1,2}\subset \mathrm{resCNN}_{r,\ell}$, значит УА сохраняется.\\
\emph{(б) Неостаточный случай.}
Любой остаток $x\mapsto x+f(x)$ можно реализовать неостаточной композицией, если разрешить больше каналов, например $r=3$
(используется тождество $\mathrm{ReLU}(x)+(-1)\mathrm{ReLU}(-x)=x$).
Это даёт УА для $\mathrm{eqCNN}_{3,2}$.

Важный дальнейший шаг --- <<доводка>> до $r=2$:

\begin{lemma}[Локальная симуляция остаточных сетей неостаточными при $r=2$.]
\label{lem:res_to_nonres}
Для любого $G\in \mathrm{eqresCNN}_{1,2}$ и компакта $K\subset X$ существует $H\in \mathrm{eqCNN}_{2,2}$, такое что
$H(x)=G(x)$ для всех $x\in K$. 
\end{lemma}
Из неё следует: раз $\mathrm{eqresCNN}_{1,2}$ эквивариантно-универсально, то и $\mathrm{eqCNN}_{2,2}$ эквивариантно-универсально,
а затем по предложению~\ref{prop:pooling_bridge} $\mathrm{CNN}_{2,2}$ становится инвариантно-универсальным.
Включение $\mathrm{CNN}_{2,2}\subset \mathrm{CNN}_{r,\ell}$ для $r\ge 2$, $\ell\ge 2$ завершает доказательство теоремы.

\begin{proof}[Доказательство Леммы \ref{lem:res_to_nonres}].
Пусть
\[
G = f_M\circ\cdots\circ f_1,\qquad f_i(x)=x+ v^i\,\sigma(w^i*x+b^i\mathbf{1}),
\]
где $\sigma=\mathrm{ReLU}$ и $\mathrm{supp}(w^i)\le 2$.
Положим $\gamma_i\coloneq f_i\circ\cdots\circ f_1$ и $\gamma_0\coloneq \mathrm{id}$.
Каждое $\gamma_i$ липшицево, а $K$ компактен, поэтому найдётся $R>0$ такое, что
$\|\gamma_i(x)\|_\infty\le R$ для всех $x\in K$ и $i=0,\dots,M$.

Определим последовательность неостаточных слоёв (каждый является элементом $\mathcal{F}_{2,2}$):
\begin{align}
u_0(x) &\coloneq  \sigma(x+R\mathbf{1}),
\label{eq:u0}\\
u_i(x) &\coloneq  \sigma(x) + v^i\,\sigma\!\Bigl(w^i*x + \bigl(b^i-(\sum_k [w^i]_k)R\bigr)\mathbf{1}\Bigr),\qquad i=1,\dots,M.
\label{eq:ui}
\end{align}
Рассмотрим композиции $\eta_s\coloneq u_s\circ\cdots\circ u_0$.
Тогда $\eta_s\in \mathrm{eqCNN}_{2,2}$.

Покажем по индукции, что
\begin{equation}
\eta_i(x)=\gamma_i(x)+R\mathbf{1}\qquad \forall x\in K,\ i=0,\dots,M.
\label{eq:eta_gamma_shift}
\end{equation}
База $i=0$ следует из \eqref{eq:u0}, так как $x+R\mathbf{1}\ge 0$ на $K$ и $\sigma$ на неотрицательных аргументах тождественна.
Шаг: пусть \eqref{eq:eta_gamma_shift} верно для $i$. Тогда для $x\in K$
\[
\eta_{i+1}(x)=u_{i+1}(\eta_i(x))
=\sigma(\gamma_i(x)+R\mathbf{1})
+ v^{i+1}\sigma\!\Bigl(w^{i+1}*(\gamma_i(x)+R\mathbf{1}) + \bigl(b^{i+1}- (\sum_k [w^{i+1}]_k)R\bigr)\mathbf{1}\Bigr).
\]
Так как $\gamma_i(x)+R\mathbf{1}\ge 0$, имеем $\sigma(\gamma_i(x)+R\mathbf{1})=\gamma_i(x)+R\mathbf{1}$.
Кроме того, $w^{i+1}*(R\mathbf{1})=(\sum_k [w^{i+1}]_k)R\mathbf{1}$.
Следовательно,
\[
\eta_{i+1}(x)=\gamma_i(x)+R\mathbf{1}+v^{i+1}\sigma(w^{i+1}*\gamma_i(x)+b^{i+1}\mathbf{1})
=\gamma_{i+1}(x)+R\mathbf{1}.
\]
Индукция завершена.

Наконец, положим $u_{M+1}(x)\coloneq \sigma(x)-\sigma(R\mathbf{1})$ и определим
\[
H(x)\coloneq u_{M+1}(\eta_M(x)).
\]
Тогда для $x\in K$ из \eqref{eq:eta_gamma_shift} получаем
$H(x)=\sigma(\gamma_M(x)+R\mathbf{1})-\sigma(R\mathbf{1})=\gamma_M(x)=G(x)$, так как $\gamma_M(x)+R\mathbf{1}\ge 0$.
Построение даёт $H\in\mathrm{eqCNN}_{2,2}$ и требуемое равенство на $K$.
\end{proof}

\underline{Доказательство пункта (2): оптимальность размера ядра.}
Пусть $\min_s \ell_s=1$ и, не теряя общности, $\ell_1=1$.
Тогда легко видеть, что любые сети из $\mathrm{CNN}_{r,\ell}+\mathbb{R}$, $\mathrm{resCNN}_{r,\ell}+\mathbb{R}$ (и даже $\mathrm{CODE}_{r,\ell}+\mathbb{R}$)
лежат в вспомогательном классе $\mathcal{H}+\mathbb{R}$, где $\mathcal{H}$ состоит из функций вида
\[
H(x)= g(\varphi(x)),
\qquad
[\varphi(x)]_{I:}=\psi(x_{I:})\ \ \text{для некоторой }\psi:X_1\to X_1,
\]
то есть каждая <<строка>> $x_{I:}$ (фиксируем первый индекс $I$) обрабатывается одинаково и независимо, а затем результаты агрегируются глобальным пулингом.

Покажем, что $\mathcal{H}+\mathbb{R}$ \emph{не} обладает УА.
Положим $K=[0,1]^n$. Определим функцию
\[
F(x)\coloneq \prod_{i_1> i_2}\Bigl(\psi([x]_{i_1:})-\psi([x]_{i_2:})\Bigr),
\qquad
\psi(y)\coloneq \prod_{i'} [y]_{i'}.
\]
Заметим: $F$ сдвиг-инвариантна (все строки участвуют симметрично, а перемножение по парам инвариантно к циклической перестановке строк).

Рассмотрим подмножества:
\[
K_1\coloneq \{x\in K:\ x_{1:}\succ x_{2:}\succ\cdots\succ x_{n_1:}\},
\qquad
K_2\coloneq \{x\in K:\ x_{2:}\succ x_{1:}\succ x_{3:}\succ\cdots\succ x_{n_1:}\},
\]
где $z_1\succ z_2$ означает $\min_i[z_1]_i\ge \max_i[z_2]_i$.
Определим преобразование $\tau$, меняющее местами первую и вторую строки:
\[
[\tau(x)]_{2:}=[x]_{1:},\qquad [\tau(x)]_{1:}=[x]_{2:},\qquad [\tau(x)]_{i:}=[x]_{i:}\ (i\ne 1,2).
\]
Тогда $\tau(K_1)=K_2$, а $\tau$ сохраняет меру. По определению $\mathcal{H}$ имеем $(H\circ\tau)(x)=H(x)$ для $x\in K_1$.
С другой стороны, из построения $F$ следует $F\circ\tau=-F$ (перестановка двух строк меняет знак произведения по парам).
Следовательно,
\[
2\|F\|_{L^p(K_1)}
=\|F-(F\circ\tau)\|_{L^p(K_1)}
\le \|H-(H\circ\tau)\|_{L^p(K_1)} + \|F-H\|_{L^p(K)} + \|F\circ\tau-H\circ\tau\|_{L^p(K)}.
\]
Первое слагаемое равно нулю, а последние два равны (так как $\tau$ сохраняет меру), поэтому
\[
2\|F\|_{L^p(K_1)}\le 2\|F-H\|_{L^p(K)}\quad\Rightarrow\quad
\|F-H\|_{L^p(K)}\ge \|F\|_{L^p(K_1)}=:\varepsilon_0>0.
\]
Значит $F$ нельзя произвольно точно приблизить элементами $\mathcal{H}+\mathbb{R}$, УА отсутствует.
Так как рассматриваемые сети содержатся в $\mathcal{H}+\mathbb{R}$, они тоже не обладают УА.

\underline{Доказательство пункта (3): почему одного канала недостаточно (неостаточный случай).}
Рассмотрим $G\in \mathrm{CNN}_{1,\infty}+\mathbb{R}$.
Тогда $G$ --- непрерывная кусочно-линейная функция на $X\cong\mathbb{R}^{|n|}$, и прямой расчёт даёт:

\emph{Наблюдение.} Существует вектор $g\in X$, такой что для почти всех $x$ (по мере Лебега) градиент $\nabla G(x)$ равен либо $0$, либо $g$.

Из наблюдения следует: для фиксированного направления $\xi\in \partial B(0,1)$ функция одной переменной
\[
\gamma(t)\coloneq G(t\xi)
\]
почти всюду имеет производную $\gamma'(t)\in\{0,\ \xi\cdot g\}$.
На полусфере $\{\xi:\ \xi\cdot g<0\}$ получаем, что $\gamma'(t)\le 0$ почти всюду, то есть $\gamma(t)$ \emph{не возрастает}.

Теперь берём целевую функцию $F(x)=\|x\|$ и компакт $B(0,1)$.
По замене переменных в полярных координатах (в $\mathbb{R}^{|n|}$):
\begin{equation}
\int_{x\in B(0,1)} |F(x)-G(x)|^p\,dx
=
\int_{\xi\in\partial B(0,1)}\int_0^1 |F(t\xi)-G(t\xi)|^p\,t^{|n|-1}\,dt\,dS(\xi).
\label{eq:polar}
\end{equation}
На полусфере $\{\xi:\xi\cdot g<0\}$ имеем $f(t)\coloneq F(t\xi)=t$ (строго возрастает), а $g(t)\coloneq G(t\xi)$ (не возрастает).
Снизу оценим вклад каждого такого луча через следующую лемму.

\begin{lemma}
\label{lem:monotone_to_constant}
Для возрастающей $f\colon[a,b]\to\mathbb{R}$ верно
\[
\inf_{\substack{g \searrow [a,b]}}\int_a^b |f-g|^p
=
\inf_{c \in \mathbb{R}}\int_a^b |f-c|^p.
\]
\end{lemma}
\begin{proof}[Доказательство леммы]
  
\end{proof}
Неравенство ``$\le$'' очевидно.
Для ``$\ge$'' возьмём произвольную невозрастающую $g$.
Положим $\tilde g(t)\coloneq g(t_0)$, где $t_0$ --- точка пересечения $f(t_0)=g(t_0)$, если она существует, иначе $\tilde g(t)\coloneq \{g(b) \text{ при } f(a) < g(a),\ g(a) \text{ при } f(a) > g(a)\}$.
Тогда для всех $t$ легко проверить $|f(t)-g(t)|\ge |f(t)-\tilde g(t)|$ (геометрически: невозрастающая кривая лежит по одну сторону от возрастающей),
откуда и следует требуемое.

Вернёмся к оценке.
Для $\xi\cdot g<0$ применяем Лемму~\ref{lem:monotone_to_constant} на отрезке $[1/2,1]$ и получаем:
\[
\int_0^1 |t-G(t\xi)|^p t^{|n|-1}dt
\ \ge\
\int_{1/2}^1 |t-G(t\xi)|^p dt\cdot \left(\tfrac12\right)^{|n|-1}
\ \ge\
\left(\tfrac12\right)^{|n|-1}\inf_{a\in[1/2,1]}\int_{1/2}^1 |t-a|^p dt.
\]
Последний инфимум вычисляется явно:
\[
\inf_{a\in[1/2,1]}\int_{1/2}^1 |t-a|^p dt
=\inf_{a\in[1/2,1]}\frac{(1-a)^{p+1}+(a-1/2)^{p+1}}{p+1}
=\frac{2\cdot (1/4)^{p+1}}{p+1}
\quad\](достигается при a=3/4)

Обозначая
\[
C_p\coloneq 2^{-|n|+1}\cdot \frac{2\cdot(1/4)^{p+1}}{p+1}>0,
\]
получаем, что для всех $\xi$ с $\xi\cdot g<0$ верно
$\int_0^1 |t-G(t\xi)|^p t^{|n|-1}dt\ge C_p$.
Интегрируя по полусфере в \eqref{eq:polar}, получаем нижнюю оценку на $L^p$-ошибку:
\[
\int_{B(0,1)} |F(x)-G(x)|^p dx \ \ge\ C_p\cdot \frac{\alpha}{2},
\]
где $\alpha$ --- мера Лебега поверхности $\partial B(0,1)$.
Следовательно, $F(x)=\|x\|$ нельзя приблизить сколь угодно точно функциями из $\mathrm{CNN}_{1,\infty}+\mathbb{R}$ на $B(0,1)$,
и УА отсутствует.
\end{proof}
\end{document}
